\chapter{Posterior computation for flexible dynamic quantile models}
\label{app:exdqlm-technical}

This appendix collects the posterior blocks that distinguish the procedures in
Chapter~\ref{ch:research-a}: the exAL scale--skewness approximation, the
transfer-function augmentation, the regularized-horseshoe update, and
posterior-predictive synthesis. Standard Gaussian filtering, smoothing, and
backward-sampling recursions follow \citet{exd-west1997bayesian} and are not
reproduced here. Distributional conventions are stated in
Subsection~\ref{exd:sec:app_distributional_conventions}.

% migration-id: MIG-CH02-01
\section{Posterior targets and state augmentation}
\label{app:exdqlm-posterior}

For the dynamic quantile linear model, the augmented AL hierarchy is
\begin{align}
y_t\mid\vect{\theta}_t,\sigma,v_t
&\sim
\N\left\{\mat{F}_t^\top\vect{\theta}_t+A(p_0)v_t,
\ \sigma B(p_0)v_t\right\},\nonumber\\
\vect{\theta}_t\mid\vect{\theta}_{t-1}
&\sim \N(\mat{G}_t\vect{\theta}_{t-1},\mat{W}_t),
\qquad
v_t\mid\sigma\sim\Exp(\text{rate}=1/\sigma).
\label{exd:eq:dyn_dqlm_hierarchy}
\end{align}
Its posterior density is proportional to the initial-state and scale priors
times the product of the observation, latent-scale, and state-evolution
densities. Conditional on \((\sigma,\vect v)\), this is a Gaussian state-space
model with observation offset \(A(p_0)v_t\) and variance
\(\sigma B(p_0)v_t\).

The exDQLM adds \(s_t\sim\TNp(0,1)\) and replaces the observation equation by
\begin{equation}
y_t\mid\vect{\theta}_t,\sigma,\gamma,v_t,s_t
\sim
\N\left\{
\mat{F}_t^\top\vect{\theta}_t+C_\gamma\sigma|\gamma|s_t+A_\gamma v_t,
\ \sigma B_\gamma v_t
\right\}.
\label{exd:eq:dyn_exdqlm_obs}
\end{equation}
Thus the joint target is
\begin{align}
p(\vect{\theta}_{0:T},\vect v,\vect s,\sigma,\gamma\mid\vect y)
&\propto p(\vect{\theta}_0)p(\sigma)\pi_\gamma(\gamma)
\prod_{t=1}^T p(y_t\mid\vect{\theta}_t,\sigma,\gamma,v_t,s_t)
\nonumber\\[-2pt]
&\quad{}\times
\prod_{t=1}^T p(v_t\mid\sigma)p(s_t)
p(\vect{\theta}_t\mid\vect{\theta}_{t-1}).
\label{exd:eq:dyn_exdqlm_joint}
\end{align}
Given the augmentation variables and scale--skewness parameters, the state
trajectory again has a Gaussian full conditional.

For transfer-function models, let
\(\mat x_t=(x_{1t},\ldots,x_{kt})^\top\). The base state is augmented by the
accumulated effect \(\zeta_t\) and instantaneous coefficients
\(\vect\psi_t\):
\begin{equation}
\vect{\theta}^{\mathrm{aug}}_t
=
(\vect{\theta}^{\mathrm{base}\top}_t,\zeta_t,\vect\psi_t^\top)^\top,
\qquad
\mat G^{tf}_t=
\begin{pmatrix}
\lambda & \mat x_t^\top\\
\mat 0_k & \mat I_k
\end{pmatrix}.
\label{exd:eq:tf_aug_state}
\end{equation}
Consequently,
\(\zeta_t=\lambda\zeta_{t-1}+\mat x_t^\top\vect\psi_{t-1}+\omega_t\),
while \(\vect\psi_t\) follows a random walk. Once this block is appended, the
same dynamic MCMC or variational calculations apply. In the static model,
\(\mat F_t^\top\vect\theta_t\) is replaced by
\(\vect x_i^\top\vect\beta\); conditional on the augmentation variables, the
coefficient update is weighted Gaussian regression.

\section{Structured scale--skewness approximation}
\label{app:exdqlm-ldvb}

The nonconjugate component of the exAL posterior is the
\((\sigma,\gamma)\) block. The dynamic variational family is
\begin{equation}
q(\vect{\theta}_{0:T},\vect v,\vect s,\sigma,\gamma)
=q(\vect{\theta}_{0:T})
\prod_{t=1}^Tq(v_t)q(s_t)\,
q_{\sigma\gamma}(\sigma,\gamma),
\label{exd:eq:dyn_exdqlm_mf}
\end{equation}
with \(q(\vect\beta)\) replacing the state factor in static regression. The
current default preserves dependence through the structured factor
\(q(\gamma)q(\sigma\mid\gamma)\). For each admissible \(\gamma\), the
conditional scale factor has a GIG kernel. Numerical quadrature on the bounded
logit coordinate
\[
z_\gamma=\logit\{(\gamma-L)/(U-L)\}
\]
approximates \(q(\gamma)\), and conditional GIG moments supply the required
scale expectations. Approximate draws first select \(\gamma\) using the
quadrature weights and then draw \(\sigma\) from the corresponding GIG factor.

The conjugate coordinates require the following scale--skewness moments:
\begin{align}
\kappa_1 &= \E\{1/(\sigma B_\gamma)\},&
\kappa_2 &= \E\{A_\gamma/(\sigma B_\gamma)\},&
\kappa_3 &= \E\{A_\gamma^2/(\sigma B_\gamma)\},\nonumber\\
\kappa_4 &= \E\{C_\gamma|\gamma|/B_\gamma\},&
\kappa_5 &= \E\{\sigma C_\gamma^2\gamma^2/B_\gamma\},&
\kappa_6 &= \E\{A_\gamma C_\gamma|\gamma|/B_\gamma\}.
\label{exd:eq:dyn_exdqlm_kappa}
\end{align}
If \(m_{1/v,t}=\E(v_t^{-1})\) and \(\bar s_t=\E(s_t)\), the Gaussian
state update uses precision and pseudo-response
\begin{equation}
\phi_t=\kappa_1m_{1/v,t},
\qquad
y_t^{\mathrm{VB}}
=\frac{y_t\phi_t-\kappa_2-\kappa_4\bar s_t m_{1/v,t}}{\phi_t}.
\label{exd:eq:dyn_exdqlm_vb_pseudo}
\end{equation}
This converts the state coordinate into a standard Gaussian filtering and
smoothing update.

Let \(m_{\eta,t}=\E(\mat F_t^\top\vect\theta_t)\),
\(S_{\eta,t}=\Var(\mat F_t^\top\vect\theta_t)\), and
\(S_{\Delta,t}=(y_t-m_{\eta,t})^2+S_{\eta,t}\). The remaining latent factors
are
\begin{align}
q(v_t)&=\GIG\left(
\tfrac12,
\kappa_1S_{\Delta,t}-2\kappa_4(y_t-m_{\eta,t})\E(s_t)
+\kappa_5\E(s_t^2),
\kappa_3+2\E(\sigma^{-1})
\right),\nonumber\\
q(s_t)&=\TNp(\mu_{s,t}^{\mathrm{VB}},V_{s,t}^{\mathrm{VB}}),
\qquad
V_{s,t}^{\mathrm{VB}}=\{1+\kappa_5\E(v_t^{-1})\}^{-1},\nonumber\\
\mu_{s,t}^{\mathrm{VB}}
&=V_{s,t}^{\mathrm{VB}}
\{\kappa_4(y_t-m_{\eta,t})\E(v_t^{-1})-\kappa_6\}.
\label{exd:eq:dyn_exdqlm_vb_vs}
\end{align}
Static regression replaces the state moments by
\(\vect x_i^\top\vect m_\beta\) and
\(\vect x_i^\top\mat\Sigma_\beta\vect x_i\).

The earlier Laplace--delta factor instead approximates transformed
\((\log\sigma,z_\gamma)\) coordinates by a Gaussian at the local mode and
uses second-order delta moments. It is retained for comparison, whereas the
structured factor is the default because it preserves the conditional scale
kernel. Both are approximations to the declared posterior target; MCMC remains
the reference when tail behavior or interval width is consequential.

\section{Regularized-horseshoe block}
\label{app:exdqlm-horseshoe}

For static regression, the \code{rhs\_ns} option uses the conditionally
conjugate regularized horseshoe of \citet{exd-nishimura2023shrunken}. If
\(\mathcal A\) indexes the coefficients subject to shrinkage, then
\begin{align}
\beta_j\mid\tau^2,\lambda_j^2,\zeta^2
&\propto
\N(\beta_j\mid0,\tau^2\lambda_j^2)
\N(\beta_j\mid0,\zeta^2),\nonumber\\
\lambda_j^2\mid\nu_j&\sim\IG(1/2,1/\nu_j),&
\nu_j&\sim\IG(1/2,1),\nonumber\\
\tau^2\mid\xi&\sim\IG(1/2,1/\xi),&
\xi&\sim\IG(1/2,1/\tau_0^2),\nonumber\\
\zeta^2&\sim\IG(a_\zeta,b_\zeta).&&
\label{exd:eq:rhs_ns_hierarchy}
\end{align}
The active prior precision is
\begin{equation}
\omega_j^{\mathrm{prior}}
=\frac{1}{\tau^2\lambda_j^2}+\frac{1}{\zeta^2}.
\label{exd:eq:rhs_ns_precision}
\end{equation}
The MCMC full conditionals for \(\lambda_j^2,\nu_j,\tau^2,\xi\), and a
random \(\zeta^2\) are inverse-gamma. Their mean-field counterparts have the
same families, with rates obtained by replacing squared coefficients and
inverse scales by their current expectations. Hence this shrinkage block adds
closed-form coordinate updates to either the static AL or exAL calculation;
it does not alter the likelihood target.

\section{Forecasting and posterior-predictive synthesis}
\label{app:exdqlm-synthesis}

At forecast origin \(\tilde t\), each posterior or variational draw propagates
the state distribution through the evolution equation. Response-level draws
then satisfy
\begin{equation}
Y_{\tilde t+k}^{\mathrm{rep},(m)}
\sim
\exAL_{p_0}\{
\mat F_{\tilde t+k}^\top\vect\theta_{\tilde t+k}^{(m)},
\sigma^{(m)},\gamma^{(m)}
\},
\label{exd:eq:forecast_postpred}
\end{equation}
with \(\gamma^{(m)}=0\) in the AL special case. Calibration plots use
one-step-ahead standardized forecast errors. CRPS is computed from
response-level predictive draws through the integrated quantile-score
identity, whereas PPLC evaluates check loss at the fitted level. These
quantities therefore assess different aspects of predictive performance.

The synthesis step combines predictive draws from separately fitted models at
ordered levels \(p_1<\cdots<p_K\). It is a post hoc predictive construction,
not a joint multi-quantile posterior.

\begin{algorithmblock}{Posterior-predictive synthesis}
\label{exd:alg:predictive-synthesis}
Extract response-level posterior-predictive draws at each fitted level; form
the quantile-specific curve for each time point and draw index; apply the
requested monotonicity correction if the fitted curve crosses; interpolate
over probability levels; and return synthesized draws and predictive
summaries on a common scale.
\end{algorithmblock}
